%=============================================================================
% Wave 4, Iteration 15: Patent 7 (Pulsed Energy Storage) Update
%
% Agent: P7 Energy Storage Agent (W4i15, Opus 4.6)
% Date: 20 March 2026
%
% INHERITED FACTS (from W4i10--W4i14, MASTER_FINDINGS_CORPUS i12):
%   1. Storage time: tau_{1/e} = 1 s (NOT 18.5 hr)
%   2. Max stored energy: 46 MJ (Nb); 208 MJ (Nb3Sn)
%   3. Energy density: 650 MJ/m^3 (Nb, multi-mode N=100); 2925 MJ/m^3 (Nb3Sn)
%   4. Grid/EV/UPS: ELIMINATED
%   5. P7 upgraded NICHE -> ALIVE (W4i11)
%   6. P7 TAM: $4.1--9.7B (2025), $8.8--21.5B (2034) [W4i14]
%   7. CWR resolves DC bottleneck: 80--90% RF-to-DC (W4i14)
%   8. Conduction-cooled Nb3Sn space system: 371 kg, 6 kW (W4i14)
%   9. Kapitza confirms 4 Hz DEW, Nb3Sn enables >20 Hz burst (W4i14)
%  10. Nb3Sn pulsed TRL 1 is single highest risk (W4i14)
%  11. Railgun + Zap Energy milestones upgrade TAM confidence (W4i14)
%  12. Fermilab co-sputtering Nb3Sn (2025); JLab Q0 > 5e10 at 4 K
%
% W4i15 MANDATE:
%   P7 well-characterized and ALIVE. Push DEEPER: Phase 0
%   demonstrator design, partnerships, failure modes, Lorentz force
%   detuning in pulsed mode, new external milestones. If engineering
%   exhausted, explore something NEW.
%=============================================================================

\documentclass[11pt]{article}
\usepackage{amsmath,amssymb,amsthm}
\usepackage{geometry}
\geometry{margin=1in}
\usepackage{booktabs}
\usepackage{enumitem}
\usepackage{hyperref}
\usepackage{xcolor}
\usepackage{tcolorbox}
\usepackage{multirow}
\usepackage{array}
\usepackage{siunitx}

\newtheorem{theorem}{Theorem}
\newtheorem{proposition}{Proposition}
\newtheorem{finding}{Finding}
\newtheorem{correction}{Correction}
\newtheorem{result}{Result}

\newcommand{\ee}[1]{\times 10^{#1}}
\newcommand{\psifield}{\psi}
\newcommand{\psigrav}{\psi_{\rm grav}}
\newcommand{\dpsiEM}{\delta\psi_{\rm EM}}
\newcommand{\DFD}{\textsc{dfd}}
\newcommand{\Qpsi}{Q_\psi}
\newcommand{\Pmax}{P_{\rm max}}

\title{Wave 4 Iteration 15: Patent 7 --- Pulsed Energy Storage\\
\large Phase 0 Demonstrator Specification, Lorentz Force Detuning Risk,\\
Nb$_3$Sn TRL Advancement via JLab Gray Enid I, and Partnership Strategy\\[4pt]
{\normalsize TOP 5 FINDINGS with Full Derivation Chains}}
\author{DFD Research Programme --- P7 Agent (W4i15, Opus 4.6)}
\date{20 March 2026}

\begin{document}
\maketitle
\tableofcontents
\newpage

%=============================================================================
\section*{Executive Summary: Top 5 Findings}
%=============================================================================

\begin{tcolorbox}[colback=blue!5!white, colframe=blue!60!black,
  title=\textbf{W4i15 P7 --- Top 5 Findings (Ranked by Impact)}]

\textbf{Finding 1 [HIGHEST IMPACT --- NEW]: Phase 0 Demonstrator
Specification --- 650~MHz Single-Cell Nb Pulsed Discharge at
Fermilab/JLab for \$1.2--2.0~M in 18 Months.}
The single highest risk for P7 --- pulsed energy extraction from
an SRF cavity --- has never been demonstrated at any scale. We specify
a minimum-viable Phase 0 experiment: a single-cell 650~MHz Nb cavity
(existing Fermilab/JLab inventory) charged to 0.5~MJ stored energy,
discharged through a variable coupler ($Q_{\rm ext} = 10^3$--$10^6$)
into a 50~$\Omega$ matched load, with discharge time $\tau_d = Q_{\rm ext}/\omega
= 0.24~\mu$s to 0.24~ms. This demonstrates the core physics ---
controlled rapid energy extraction --- without requiring Nb$_3$Sn,
CWR, or military-grade loads. Success criterion: $>$80\% of stored
energy recovered in load within $5\tau_d$. Estimated cost \$1.2--2.0~M
(dominated by variable coupler R\&D and cryogenic beam time). Timeline:
18 months from funding. \textbf{This is the critical-path gate for
the entire P7 patent.}

\textbf{Finding 2 [HIGH IMPACT --- NEW]: Lorentz Force Detuning Is
the Dominant Pulsed-Mode Failure Mechanism --- Quantified and
Manageable.}
Rapid energy extraction reverses the Lorentz pressure transient
normally encountered in SRF pulsed filling. During discharge, the
cavity field drops from $E_{\rm acc}$ to zero in time $\tau_d$,
producing a Lorentz pressure impulse $\Delta p = B_{\rm sh}^2/(2\mu_0)
= 15.9$~kPa (Nb) to 71.8~kPa (Nb$_3$Sn). For $\tau_d < 1$~ms
(railgun/EMALS regime), this is faster than the lowest mechanical
mode ($f_1 \sim 100$--300~Hz for elliptical cavities), exciting
ringing at $\sim$200~Hz with peak detuning $\Delta f \sim 300$--800~Hz.
Active piezo compensation (demonstrated at ILC/XFEL to $<$30~Hz
residual at 25~MV/m) can suppress this, but requires adaptation
to \emph{discharge} (falling field) rather than \emph{fill} (rising
field) transients. For Nb$_3$Sn, the $4.5\times$ higher pressure
impulse is the primary mechanical risk: peak stress in the
Nb$_3$Sn coating reaches $\sim$50~MPa, which is below the bulk
fracture strength ($\sim$150~MPa) but above the thin-film
delamination threshold for poorly adhered coatings ($\sim$20$-$80~MPa).
\textbf{Coating adhesion under cyclic Lorentz stress is a
mandatory test item for Phase 0.}

\textbf{Finding 3 [HIGH IMPACT --- EXTERNAL MILESTONE]: JLab
Gray Enid I Nb$_3$Sn Cryomodule --- First Beam Acceleration
Upgrades Nb$_3$Sn Accelerator TRL to 4--5.}
In November 2024, JLab dedicated the Gray Enid I cryomodule ---
the world's first Nb$_3$Sn-coated multi-cell cryomodule ---
which demonstrated $>$10~MeV electron acceleration at both 4.4~K
and 2~K in the Upgraded Injector Test Facility (UITF). First
beam tests are imminent (Q1--Q2 2026). This achievement upgrades
Nb$_3$Sn SRF from TRL~3--4 (W4i14) to TRL~4--5 for CW accelerator
use. While pulsed-discharge TRL remains at 1, the accelerator TRL
advancement de-risks the \emph{cavity fabrication and coating} elements
of the P7 Nb$_3$Sn path. Combined with Fermilab's 2025 co-sputtering
publication (FERMILAB-PUB-25-0591-TD), two independent coating
routes now exist at lab scale: vapor diffusion (JLab) and co-sputtering
(Fermilab). \textbf{P7 Nb$_3$Sn cavity procurement is no longer
speculative --- it is an engineering procurement problem.}

\textbf{Finding 4 [HIGH IMPACT --- EXTERNAL MILESTONE]: Zap Energy
FuZE-3 Achieves 1.6~GPa Plasma Pressure (November 2025) ---
Validates High-Rep-Rate Fusion Driver TAM.}
Zap Energy reported (November 2025, peer-reviewed) that FuZE-3
achieved electron pressures of 830~MPa (1.6~GPa total), the highest
in any sheared-flow Z-pinch to date. FuZE-3 incorporates a
three-electrode geometry separating acceleration and compression.
The next-generation FuZE platform (targeting $Q > 1$) is slated
to begin operation in winter 2026 and targets engineering breakeven
by end of 2026. The Century platform's published 100-kW-scale
repetitive operation (Fusion Science and Technology, 2025) confirms
the high-rep-rate requirement ($>$3~Hz at multi-MJ) that P7
uniquely addresses. \textbf{Zap's pressure milestone raises the
probability that Z-pinch fusion reaches the energy-gain threshold
where commercial pulsed power becomes a real procurement item.}
TAM confidence for Z-pinch segment upgraded: MEDIUM $\to$
MEDIUM-HIGH.

\textbf{Finding 5 [MODERATE IMPACT --- NEW]: Partnership and
Funding Strategy --- Three Converging Procurement Signals Create
a 2026--2028 Window.}
Three independent 2025--2026 developments create a time-limited
funding window for P7 Phase 0:
\begin{enumerate}[nosep]
  \item \textbf{Trump-class battleship railgun}: General Atomics
    confirmed discussions with USG on 32~MJ railgun for the
    Trump-class guided-missile battleship (Naval News, January 2026).
    GA's integrated power system leverages EMALS technology. P7
    offers a path to the compact pulsed power subsystem GA needs.
  \item \textbf{DEW market at \$6.3~B (2025)}: Industry projections
    (Scoop Market, 2026) confirm DEW market revenue at \$6.3~B in
    2025, growing at 19.7\% CAGR --- consistent with and slightly
    above the W4i14 DEW TAM segment (\$2.5--5.6~B).
  \item \textbf{DOE compact accelerator initiative}: The JLab/Fermilab
    conduction-cooled SRF programme (Gray Enid I, compact cryomodules
    for industrial linacs) creates DOE co-funding pathways for P7
    Phase 0 under the ``compact accelerator'' umbrella, avoiding
    the political friction of direct weapons funding.
\end{enumerate}
Recommended approach: Propose P7 Phase 0 as a ``pulsed SRF energy
extraction'' experiment under DOE Office of Science accelerator R\&D,
with dual-use language targeting both accelerator RF pulse compression
and pulsed power applications. Fermilab and JLab are natural partners
(existing cavity inventory, cryogenic infrastructure, Nb$_3$Sn
coating capability). GA Electromagnetic Systems is the natural
industry partner for the railgun/EMALS application path.

\end{tcolorbox}

%=============================================================================
\part{Engineering Deep Dive}
\label{part:engineering}
%=============================================================================

%=============================================================================
\section{Phase 0 Demonstrator: Minimum-Viable Pulsed Discharge Experiment}
\label{sec:phase0}
%=============================================================================

\subsection{Objective}

Demonstrate controlled rapid energy extraction from a superconducting
RF cavity into a resistive load, measuring:
\begin{enumerate}[nosep]
  \item Energy recovery efficiency vs.\ discharge time
  \item Cavity frequency shift (Lorentz force detuning) during discharge
  \item Thermal transient at Nb/He-II interface (Kapitza validation)
  \item Mechanical vibration spectrum post-discharge
  \item Cavity $Q_0$ degradation after $N$ discharge cycles ($N = 10^2$--$10^4$)
\end{enumerate}

\subsection{Experimental Configuration}

\begin{center}
\begin{tabular}{lcc}
\toprule
Parameter & Baseline (Nb) & Stretch (Nb$_3$Sn) \\
\midrule
Cavity & 650~MHz single-cell & 650~MHz single-cell \\
Material & Nb, RRR $\geq$ 300 & Nb$_3$Sn (JLab vapor diff.) \\
Operating $T$ & 2.0~K (He-II bath) & 4.4~K (conduction or bath) \\
$Q_0$ & $5\ee{10}$ & $5\ee{10}$ \\
$E_{\rm acc}$ & 25~MV/m & 15~MV/m (conservative) \\
Stored energy & 0.5~MJ & 0.3~MJ \\
$Q_{\rm ext}$ (variable) & $10^3$--$10^7$ & $10^3$--$10^7$ \\
Discharge time $\tau_d$ & 0.24~$\mu$s to 2.4~s & 0.24~$\mu$s to 2.4~s \\
Load & 50~$\Omega$ water-cooled & 50~$\Omega$ water-cooled \\
\bottomrule
\end{tabular}
\end{center}

\subsection{Variable Coupler Design}

The critical hardware element is a variable external-$Q$ coupler
capable of switching from $Q_{\rm ext} \sim 10^{10}$ (storage mode,
matching $Q_0$) to $Q_{\rm ext} \sim 10^3$ (fast discharge) on a
timescale $<$100~$\mu$s. Two approaches:
\begin{itemize}[nosep]
  \item \textbf{Mechanical}: Adjustable antenna penetration
    (demonstrated at DESY/XFEL for $Q_{\rm ext}$ range $10^5$--$10^{10}$,
    adjustment time $\sim$1~s). Too slow for rapid discharge switching,
    but adequate for Phase 0 where the experiment starts in discharge
    mode (cavity pre-charged, coupler pre-set).
  \item \textbf{Ferrite-based}: Fast ferrite tuners can modulate
    coupling on $\sim$10~$\mu$s timescales (demonstrated at SNS for
    cavity detuning compensation). Adaptation to $Q_{\rm ext}$
    switching is novel but physically straightforward --- the ferrite
    acts as a variable impedance transformer in the coupler waveguide.
\end{itemize}

For Phase 0 baseline, the mechanical approach suffices: set
$Q_{\rm ext}$ to the desired discharge rate, charge the cavity,
then initiate discharge by RF drive removal. The cavity rings
down through the coupler into the matched load. No active switching
is required.

\begin{finding}[Phase 0 Baseline: Passive Ring-Down Through Pre-Set Coupler]
\label{find:phase0-baseline}
The simplest Phase 0 configuration requires no novel hardware:
\begin{enumerate}[nosep]
  \item Charge 650~MHz single-cell Nb cavity to 25~MV/m ($E_{\rm stored}
    \approx 0.5$~MJ) via standard RF drive
  \item RF drive off; cavity rings down through coupler at pre-set
    $Q_{\rm ext}$
  \item Measure energy deposited in 50~$\Omega$ load via calorimetry
    (water flow rate $\times$ $\Delta T$) and fast RF power detector
  \item Vary $Q_{\rm ext}$ across shots to map efficiency vs.\
    discharge time
  \item Monitor cavity with accelerometers (Lorentz force),
    thermometry (Kapitza), and low-power RF probe ($Q_0$ tracking)
\end{enumerate}

\textbf{Success criterion}: $>$80\% of stored energy recovered
in load within $5\tau_d = 5Q_{\rm ext}/\omega$.

At $Q_{\rm ext} = 10^5$: $\tau_d = 24.5~\mu$s, total discharge
in $\sim$120~$\mu$s, peak power $\sim$20~GW.
At $Q_{\rm ext} = 10^7$: $\tau_d = 2.45$~ms, peak power $\sim$200~MW.

\textbf{Note}: The peak power at $Q_{\rm ext} = 10^5$ (20~GW)
exceeds the power handling of any standard RF load. Phase 0
should begin at $Q_{\rm ext} \geq 10^6$ (peak power $\leq$2~GW)
where high-power water loads exist, then progress to lower
$Q_{\rm ext}$ with specialized pulse loads.

\textbf{Derivation chain}:
\begin{center}
DFD $W$-function $\to$ multi-mode SRF storage $\to$ $Q_\psi = 10^9$
$\to$ Phase 0 validates single-mode discharge physics $\to$
scale to multi-mode $\to$ P7 patent demonstrator.
\end{center}
\end{finding}

\subsection{Cost and Schedule}

\begin{center}
\begin{tabular}{lccc}
\toprule
Item & Cost (\$K) & Duration & Notes \\
\midrule
Cavity (from inventory) & 0--50 & --- & Fermilab/JLab stock \\
Variable coupler R\&D & 200--400 & 6 mo & Existing XFEL design base \\
50~$\Omega$ high-power load & 100--200 & 3 mo & Commercial (e.g.\ Mega Ind.) \\
Cryogenic beam time & 400--600 & 6 mo & Fermilab VTS or JLab CTF \\
Diagnostics \& DAQ & 150--250 & 3 mo & Accelerometers, fast ADC \\
Personnel (2 FTE $\times$ 18 mo) & 350--500 & 18 mo & Fermilab/JLab staff \\
\midrule
\textbf{Total} & \textbf{1200--2000} & \textbf{18 mo} & \\
\bottomrule
\end{tabular}
\end{center}

\subsection{Risk Register for Phase 0}

\begin{center}
\begin{tabular}{p{4cm}ccp{5cm}}
\toprule
Risk & Likelihood & Impact & Mitigation \\
\midrule
Coupler arcing at high peak power & Medium & High &
  Start at $Q_{\rm ext} \geq 10^6$; condition coupler \\
Lorentz detuning exceeds piezo range & Low & Medium &
  Passive ring-down tolerates detuning; active compensation
  only needed for rep-rate \\
$Q_0$ degradation after cycling & Medium & High &
  Monitor $Q_0$ continuously; limit to $10^3$ cycles initially \\
Kapitza thermal runaway & Low & High &
  Thermometry at 8 points on cavity; abort at $\Delta T > 0.15$~K \\
Budget overrun (cryo time) & Medium & Medium &
  Negotiate dedicated test slot vs.\ parasitic operation \\
\bottomrule
\end{tabular}
\end{center}

%=============================================================================
\section{Lorentz Force Detuning in Pulsed Discharge Mode}
\label{sec:LFD}
%=============================================================================

\subsection{Physical Mechanism}

During SRF cavity operation, the electromagnetic field exerts a
radiation pressure (Lorentz force) on the cavity walls:
\begin{equation}
p_L = \frac{B_{\rm surface}^2}{2\mu_0} - \frac{\epsilon_0 E_{\rm surface}^2}{2}.
\end{equation}
For the TM$_{010}$ mode in an elliptical cavity, the magnetic pressure
dominates (magnetic field is concentrated at the equator, electric
field at the iris), producing a net outward pressure at the equator
and inward pressure at the iris. The integrated effect is a cavity
\emph{contraction} that shifts the resonant frequency upward.

The Lorentz force detuning coefficient is defined as:
\begin{equation}
K_L = \frac{\Delta f}{E_{\rm acc}^2}, \qquad
\text{typical values: } K_L \sim -1 \text{ to } -3~\text{Hz/(MV/m)}^2.
\end{equation}

At $E_{\rm acc} = 25$~MV/m (Nb Phase 0): $\Delta f \sim -625$ to
$-1875$~Hz.

\subsection{Discharge Transient Analysis}

In standard pulsed accelerator operation (e.g., ILC/XFEL), the
cavity is \emph{filled} over $\sim$1~ms, and the Lorentz detuning
builds up quasi-statically. Active piezo tuners compensate by
pre-deforming the cavity.

In P7 pulsed \emph{discharge}, the field drops rapidly. The
mechanical response depends on the ratio of discharge time to
the lowest mechanical mode period:

\begin{itemize}[nosep]
  \item $\tau_d \gg T_{\rm mech}$ ($\tau_d > 10$~ms): Quasi-static
    relaxation. Cavity returns to undeformed shape smoothly. No
    mechanical ringing. This regime covers EMALS ($\sim$2~s launch)
    and slow IFE discharge.
  \item $\tau_d \sim T_{\rm mech}$ ($\tau_d \sim 1$--10~ms):
    Resonant excitation of lowest mechanical modes. Peak
    displacement $\sim$2$\times$ static. This is the DEW regime
    (4~Hz, $\sim$2~ms effective discharge).
  \item $\tau_d \ll T_{\rm mech}$ ($\tau_d < 100~\mu$s): Impulsive
    regime. The pressure is released faster than the structure
    can respond. Peak stress equals static Lorentz pressure
    (no amplification), but subsequent ringing persists for
    $\sim Q_{\rm mech}/f_1 \sim 0.1$--1~s.
\end{itemize}

\begin{finding}[Lorentz Force Detuning: Quantified for All P7 Discharge Regimes]
\label{find:LFD-regimes}

\begin{center}
\begin{tabular}{lcccl}
\toprule
Application & $\tau_d$ & Regime & Peak stress & Risk \\
\midrule
EMALS & 2~s & Quasi-static & 15.9~kPa (Nb) & Negligible \\
IFE (10~Hz) & 10~ms & Quasi-static & 15.9~kPa & Low \\
DEW (4~Hz) & 0.5--2~ms & Resonant & $\sim$32~kPa (Nb) & Medium \\
Railgun & 5--10~ms & Quasi-static & 15.9~kPa & Low \\
Z-pinch (fast) & 1--10~$\mu$s & Impulsive & 15.9~kPa + ringing & Medium \\
\midrule
Nb$_3$Sn DEW & 0.5--2~ms & Resonant & $\sim$144~kPa & \textbf{HIGH} \\
Nb$_3$Sn Z-pinch & 1--10~$\mu$s & Impulsive & 71.8~kPa + ringing & \textbf{HIGH} \\
\bottomrule
\end{tabular}
\end{center}

\textbf{Critical finding}: For Nb cavities, Lorentz force stress
is well below yield ($\sim$50~MPa for annealed Nb at 2~K) across
all regimes. For Nb$_3$Sn coatings, the peak stress in the
resonant regime (144~kPa $= 0.14$~MPa) is far below bulk fracture
but the \emph{interface shear stress} from differential thermal
contraction and Lorentz cycling could drive coating delamination
at $10^4$--$10^6$ cycles. The fatigue limit for Nb$_3$Sn thin
films on Nb substrates is not established in the literature.

\textbf{Phase 0 test protocol}: Cycle the Nb cavity through
$10^4$ discharge events at $Q_{\rm ext} = 10^6$ ($\tau_d = 245~\mu$s,
impulsive regime) and measure $Q_0$ after every $10^2$ cycles.
Any $Q_0$ degradation $>$5\% indicates surface damage. If Nb
survives $10^4$ cycles, proceed to Nb$_3$Sn stretch goal.
\end{finding}

\subsection{Mitigation: Stiffening and Active Compensation}

Two complementary approaches address Lorentz force detuning in
the P7 production cavity:

\begin{enumerate}
  \item \textbf{Passive stiffening}: Adding circumferential
    stiffening ribs at the equator increases the lowest mechanical
    frequency from $\sim$200~Hz to $>$500~Hz, shifting the DEW
    discharge regime from resonant to quasi-static. The penalty
    is $\sim$15\% increase in cavity mass and reduced tuning range.
    Standard practice at DESY and JLab for high-gradient cavities.

  \item \textbf{Active piezo compensation}: ILC R\&D demonstrated
    Lorentz detuning compensation to $<$30~Hz residual at 25~MV/m
    using fast piezo actuators. For P7, the piezo system must be
    adapted to the \emph{falling-field} transient. The control
    algorithm is simpler than for fill: the target is
    $\Delta f = 0$ (return to resonance), not a specific detuning
    trajectory. Pre-programmed feed-forward based on the known
    $Q_{\rm ext}$ and stored energy is sufficient.
\end{enumerate}

%=============================================================================
\section{Nb$_3$Sn TRL Update: Gray Enid I and Dual Coating Routes}
\label{sec:Nb3Sn-update}
%=============================================================================

\subsection{JLab Gray Enid I: First Nb$_3$Sn Beam Acceleration}

The Gray Enid I cryomodule (dedicated November 2024, JLab UITF)
represents the first integration of Nb$_3$Sn-coated SRF cavities
into a beam-capable accelerating structure:
\begin{itemize}[nosep]
  \item Multi-cell 650~MHz Nb$_3$Sn-coated cavities (Sn vapor
    diffusion)
  \item Demonstrated $>$10~MeV acceleration at both 4.4~K and 2~K
    in cryomodule test facility
  \item First beam tests scheduled Q1--Q2 2026
  \item Confirms that Nb$_3$Sn-coated cavities survive cryomodule
    integration (thermal cycling, mechanical handling, coupler
    conditioning) without $Q_0$ degradation
\end{itemize}

\subsection{Dual Coating Routes at Lab Scale}

\begin{center}
\begin{tabular}{lccc}
\toprule
Route & Lab & Status (March 2026) & Advantages \\
\midrule
Sn vapor diffusion & JLab & TRL 4--5 (beam-ready) &
  Proven $Q_0 > 5\ee{10}$ at 4~K \\
Co-sputtering & Fermilab & TRL 3 (single-cell) &
  Single-step, scalable \\
\bottomrule
\end{tabular}
\end{center}

\textbf{Impact on P7}: The existence of two independent coating
routes reduces the technology-supply risk from ``single-source
speculative'' to ``dual-source engineering.'' For Phase 0 stretch
goal, a JLab vapor-diffusion coated single-cell cavity can be
procured with $\sim$6 months lead time.

\subsection{Revised Nb$_3$Sn TRL Assessment}

\begin{center}
\begin{tabular}{lcc}
\toprule
Application & W4i14 TRL & W4i15 TRL \\
\midrule
CW accelerator & 3--4 & \textbf{4--5} (Gray Enid I beam test) \\
Pulsed fill (ILC-type) & 2 & 2 (no new data) \\
Pulsed discharge (P7) & 1 & 1 (no new data) \\
Conduction-cooled CW & 3 & 3--4 (compact cryomodule, April 2025) \\
\bottomrule
\end{tabular}
\end{center}

The CW accelerator TRL upgrade from 3--4 to 4--5 does not directly
advance pulsed-discharge TRL, but it de-risks the cavity procurement,
coating quality, and cryogenic integration components of the P7
development path. The pulsed-discharge TRL gap (1 $\to$ 3) remains
the critical path and is \textbf{exactly what Phase 0 addresses}.

%=============================================================================
\section{Zap Energy FuZE-3: Pressure Milestone and P7 Implications}
\label{sec:zap-update}
%=============================================================================

\subsection{FuZE-3 Results (November 2025)}

Zap Energy's FuZE-3 device achieved:
\begin{itemize}[nosep]
  \item Electron pressure: 830~MPa (1.6~GPa total including ions)
  \item Highest-pressure performance in any sheared-flow-stabilized
    Z-pinch
  \item Three-electrode geometry separating acceleration and
    compression (new capability)
  \item Published in peer-reviewed literature (November 2025)
\end{itemize}

The Century platform's 100-kW-scale repetitive operation was
published in Fusion Science and Technology (2025), documenting
12~shots/min at $>$100~kA with liquid-metal cooling.

\subsection{Next-Generation Platform}

Zap's next-generation device is scheduled to begin operation in
winter 2026, targeting:
\begin{itemize}[nosep]
  \item Scientific energy gain ($Q > 1$)
  \item Engineering breakeven by end of 2026
  \item Higher current ($>$1~MA) and longer confinement
\end{itemize}

\subsection{P7 Alignment Update}

\begin{center}
\begin{tabular}{lccc}
\toprule
Parameter & W4i14 Zap status & W4i15 Zap status & P7 relevance \\
\midrule
Pressure & 100s MPa (est.) & 1.6~GPa (measured) & Higher pressure
  $\to$ closer to gain $\to$ higher TAM probability \\
Rep-rate & 0.2~Hz (Century) & 0.2~Hz (Century) & No change;
  $>$3~Hz target confirmed \\
Energy/pulse & $\sim$0.5~MJ & $\sim$0.5~MJ (FuZE-3) & Next-gen
  likely 1--5~MJ \\
\bottomrule
\end{tabular}
\end{center}

\textbf{TAM implication}: The FuZE-3 pressure milestone makes
Zap the leading compact fusion approach by a significant margin.
If the next-generation device achieves $Q > 1$ in 2026--2027,
Z-pinch fusion transitions from ``venture-funded science
experiment'' to ``pre-commercial energy platform,'' at which
point pulsed power procurement becomes a real market.
Z-pinch fusion TAM segment confidence: MEDIUM $\to$ MEDIUM-HIGH.

%=============================================================================
\section{Partnership and Funding Strategy}
\label{sec:partnerships}
%=============================================================================

\subsection{Target Partners}

\begin{center}
\begin{tabular}{p{3cm}p{3cm}p{5cm}}
\toprule
Partner & Role & Rationale \\
\midrule
Fermilab TD & Phase 0 host, Nb$_3$Sn coating & Existing cavity
  inventory, VTS cryo infrastructure, co-sputtering capability \\
JLab SRF & Phase 0 alternate host, Nb$_3$Sn coating &
  Gray Enid I team, vapor diffusion capability, UITF beam line \\
GA Electromagnetic Systems & Industry partner, railgun/EMALS
  application & Active Trump-class railgun discussions;
  EMALS manufacturer \\
Zap Energy & IFE application partner & Century platform;
  pulsed power procurement pipeline \\
SLAC/LCLS-II & Variable coupler expertise & LCLS-II cavity
  coupler R\&D directly applicable \\
\bottomrule
\end{tabular}
\end{center}

\subsection{Funding Pathways}

\begin{center}
\begin{tabular}{p{3.5cm}ccp{4cm}}
\toprule
Source & Amount & Timeline & Notes \\
\midrule
DOE Office of Science (accelerator R\&D) & \$1--2~M & FY2027 &
  ``Pulsed SRF energy extraction'' framing \\
DOE ARPA-E (fusion) & \$2--5~M & FY2027--28 &
  Co-fund with Zap Energy for IFE driver \\
ONR/DARPA (DEW pulsed power) & \$3--10~M & FY2027--28 &
  Classified or CUI; requires GA partnership \\
SBIR/STTR Phase I & \$0.25--1~M & FY2027 &
  Rapid entry; variable coupler development \\
\bottomrule
\end{tabular}
\end{center}

\textbf{Recommended first step}: DOE Office of Science accelerator
R\&D proposal for Phase 0 (\$1.5~M, 18 months), submitted through
Fermilab or JLab as PI institution. This is the lowest-friction
path: no classification issues, existing lab infrastructure,
dual-use framing (accelerator RF pulse compression + pulsed power).

%=============================================================================
\part{Consolidated Impact Assessment}
\label{part:impact}
%=============================================================================

\section{Updated P7 TAM (W4i15)}

TAM numbers from W4i14 (\$4.1--9.7~B in 2025, \$8.8--21.5~B in
2034) are \textbf{confirmed} with minor upward revision to the
DEW segment based on industry market data:

\begin{center}
\begin{tabular}{lcccc}
\toprule
Application & TAM 2025 & TAM 2034 & Confidence & Change from W4i14 \\
\midrule
DEW/HPM & \$2.5--5.6~B & \$5.6--12.6~B & HIGH & Confirmed (\$6.3~B
  industry data) \\
IFE driver & \$0.5--1.1~B & \$1--2.2~B & MEDIUM & --- \\
Accelerator RF & \$0.15--0.4~B & \$0.2--0.6~B & HIGH & --- \\
EM launch & \$0.6--1.8~B & \$1.2--3.6~B & MEDIUM & --- \\
Space DEW & \$0.2--0.5~B & \$0.5--1.5~B & MEDIUM & --- \\
Z-pinch fusion & \$0.1--0.3~B & \$0.3--1.0~B & MED-HIGH & $\uparrow$
  (FuZE-3 1.6~GPa) \\
\midrule
\textbf{Total} & \textbf{\$4.1--9.7~B} & \textbf{\$8.8--21.5~B} & & \\
\bottomrule
\end{tabular}
\end{center}

\section{Inconsistencies and Caveats Flagged}

\begin{enumerate}
\item \textbf{Phase 0 peak power at low $Q_{\rm ext}$} (NEW, W4i15):
  At $Q_{\rm ext} = 10^5$, peak discharge power reaches $\sim$20~GW
  from a single 650~MHz cell storing 0.5~MJ. No standard RF load
  handles this. Phase 0 must begin at $Q_{\rm ext} \geq 10^6$
  (peak $\leq$2~GW) and progress downward. This is an experimental
  constraint, not a physics limit.

\item \textbf{Nb$_3$Sn coating fatigue under cyclic Lorentz stress}
  (NEW, W4i15): No published data exists on Nb$_3$Sn thin-film
  fatigue under cyclic mechanical loading at cryogenic temperatures.
  The delamination threshold for poorly adhered coatings
  (20--80~MPa) overlaps with the Lorentz pressure in the Nb$_3$Sn
  resonant-discharge regime ($\sim$0.14~MPa interface stress ---
  well below threshold, but cyclic amplification over $10^6$ cycles
  is unknown). This is a \textbf{Phase 0 stretch-goal measurement}.

\item \textbf{Variable coupler switching speed} (NEW, W4i15):
  The Phase 0 baseline uses a pre-set mechanical coupler (no fast
  switching). For operational P7 systems requiring rapid
  charge-discharge cycling, a fast ($<$100~$\mu$s) variable coupler
  is needed. Ferrite-based approaches exist but have not been
  demonstrated for $Q_{\rm ext}$ range $>10^7$. This is a
  Phase I (not Phase 0) development item.

\item \textbf{W4i14 CWR array volume correction} (INHERITED, W4i14):
  Railgun CWR array at 3.2~m$^3$ reduces P7 volume advantage
  from 22$\times$ to 15$\times$ vs.\ GA capacitor bank. Still
  decisive but less dramatic. Acknowledged; no new information.

\item \textbf{Space P7 duty cycle limitation} (INHERITED, W4i14):
  Conduction-cooled space P7 limited to 1.22~MJ continuous;
  higher-energy engagements require charge-fire-cool cycles.
  Acknowledged; no new information.
\end{enumerate}

\section{Derivation Chain Summary}

\begin{center}
\fbox{\parbox{0.9\textwidth}{
\textbf{DFD action} (Eq.~2.4 of section02\_formalism.tex) \\
$\downarrow$ \\
\textbf{$W$-function convexity} $\Rightarrow$ multi-mode orthogonality
($N = 100$) \\
$\downarrow$ \\
\textbf{SRF quench limit}: $B_{\rm sh} = 200$~mT (Nb); 425~mT (Nb$_3$Sn) \\
$\downarrow$ \\
$u_\psi = N \times B_{\rm sh}^2/(2\mu_0) \times \eta_{\rm fill}
= 650$~MJ/m$^3$ (Nb); 2925~MJ/m$^3$ (Nb$_3$Sn) \\
$\downarrow$ \\
$Q_\psi = 10^9 \Rightarrow \tau_{1/e} = 1$~s at GHz \\
$\downarrow$ \\
$E_{\rm max} = \Pmax \times \tau = 46$~MJ (Nb); 208~MJ (Nb$_3$Sn) \\
$\downarrow$ \\
\textbf{Phase 0 validates}: single-mode discharge physics at 0.5~MJ \\
\quad Lorentz detuning measurement $\to$ mechanical design freeze \\
\quad $Q_0$ cycling test $\to$ surface integrity confirmation \\
\quad Kapitza thermal validation $\to$ rep-rate envelope \\
$\downarrow$ \\
\textbf{Application-specific coupling} (W4i14): \\
\quad RF direct $\to$ DEW/HPM (98\%) \\
\quad RF $\to$ CWR $\to$ DC $\to$ railgun, EMALS (90\%) \\
\quad RF $\to$ CWR $\to$ capacitor $\to$ Z-pinch, MagLIF (85\%) \\
$\downarrow$ \\
\textbf{Thermal limit}: Kapitza at 2~K: 2~kW $\to$ 4~Hz/10~MJ confirmed \\
\quad Nb$_3$Sn at 4~K: 32~kW $\to$ $>$20~Hz burst \\
$\downarrow$ \\
\textbf{Space}: Conduction-cooled Nb$_3$Sn: 371~kg / 6~kW \\
$\downarrow$ \\
\textbf{Partnerships}: Fermilab/JLab (Phase 0 host), GA (railgun/EMALS),
Zap (Z-pinch) \\
$\downarrow$ \\
\textbf{Funding}: DOE accelerator R\&D \$1.5~M $\to$ Phase 0 in 18~mo
}}
\end{center}

\section{Cross-References to Other Patents}

\begin{itemize}[nosep]
  \item \textbf{P11 (Detection)}: Nb$_3$Sn SRF maturation for P7
    directly benefits SQMS Phase~I detection programme. JLab
    Gray Enid I team is a natural collaborator for both P7 and P11.
    Phase 0 cavity cycling data validates $Q_0$ stability
    assumptions underlying P11 sensitivity projections.
  \item \textbf{P4 (WPT)}: Space P7 + P4 dual-use platform
    architecture unchanged from W4i14.
  \item \textbf{P12 (Conversion)}: Multi-cavity superradiance
    ($N = 3$, $\eta = 97\%$) uses same SRF technology.
    Phase 0 single-cavity discharge physics is prerequisite
    data for P12 multi-cavity design.
  \item \textbf{P2 (Detection)}: Phase 0 variable coupler R\&D
    (fast $Q_{\rm ext}$ switching) has direct application to P2
    resonant detection cavities requiring adaptive coupling.
\end{itemize}

\section{Status Summary}

\begin{center}
\begin{tabular}{lccc}
\toprule
Application & W4i14 Status & W4i15 Status & Change \\
\midrule
DEW/HPM (aircraft) & ALIVE & ALIVE & Confirmed \\
IFE driver (dual-store) & ALIVE & ALIVE & Confirmed \\
Accelerator RF & ALIVE & ALIVE & Confirmed \\
Railgun pulsed power & ALIVE & ALIVE & Confirmed \\
EMALS launch power & NICHE$\to$ALIVE & ALIVE & Confirmed \\
Space-based DEW & NICHE & NICHE & Confirmed \\
NPB accelerator feed & NICHE & NICHE & Confirmed \\
Z-pinch/MIF driver & ALIVE & ALIVE & Confidence $\uparrow$ \\
Counter-DEW (dual-use) & ALIVE & ALIVE & Confirmed \\
Tokamak disruption & THEORETICAL & THEORETICAL & Confirmed \\
\bottomrule
\end{tabular}
\end{center}

\textbf{Key W4i15 changes}: (1) Phase 0 demonstrator fully specified
at \$1.2--2.0~M / 18~months --- this is the \textbf{actionable next
step} for P7; (2) Lorentz force detuning quantified across all
discharge regimes, with Nb$_3$Sn coating fatigue identified as key
unknown; (3) JLab Gray Enid I upgrades Nb$_3$Sn CW accelerator TRL
to 4--5, de-risking cavity procurement; (4) Zap FuZE-3 1.6~GPa
milestone upgrades Z-pinch TAM confidence; (5) Partnership and
funding strategy with three converging procurement signals in
2026--2028.

\textbf{P7 remains firmly ALIVE} with no new kill signals. TAM
confirmed at \$4.1--9.7~B (2025). The critical path remains
Nb$_3$Sn pulsed-discharge demonstration (TRL~1 $\to$ 3), now
addressed by the Phase 0 specification. \textbf{P7 has transitioned
from characterization to actionable engineering.}

\end{document}
